\documentclass[12pt,a4paper,oneside]{article}
\usepackage[brazil]{babel}
\usepackage[utf8]{inputenc}
\usepackage{olymp}

\setcounter{problem}{3}

\begin{document}

\begin{problem}{Equação de segundo grau}{equacao2}{2}
 
Faça um programa que calcule ambas as raízes de uma equação de segundo grau.
%\Notes
%
%Observações, se tiver.

\InputFile

A entrada terá apenas um caso de teste.

O caso de teste é composto por três valores reais, $A$, $B$ e $C$, representando os coeficientes da equação de segundo grau.

\Notes
Neste problema, é garantido que sempre haverá duas raízes reais.


\OutputFile

Seu programa deve gerar apenas uma linha de saída, contendo as raízes da equação de segundo grau de coeficientes $A$, $B$, $C$. Esses valores devem ser impressos como números reais com precisão de 4 casas decimais, separados por um espaço em branco, seguindo a ordem mostrada nos exemplos.

%O primeiro valor é a raiz  $\frac{-b+\sqrt{\Delta}}{2a}$ e o segundo a raiz $\frac{-b-\sqrt{\Delta}}{2a}$, sendo $\Delta = b^2-4ac$.

Não se esqueça de quebrar a linha após a impressão dos valores.

\Example

\begin{example}
\exmp{
1 -2 -3
}{
3.0000 -1.0000
}%
\end{example}

\begin{example}
\exmp{
1 54 55
}{
-1.0385 -52.9615
}%
\end{example}

\begin{example}
\exmp{
1 1 -20
}{
4.0000 -5.0000
}%
\end{example}


%Comentários sobre o exemplo, se necessário.

\end{problem}

\end{document}
